% Einheit 2 – Extrempunkte
\setcounter{aufgabe}{14}
\einheitenkopf[e2][2 Extrempunkte]{Einheit 2 von 5 · Extrempunkte}
\zweigzeile{Hier lernst du, Hoch-, Tief- und Sattelpunkte zu berechnen und an einer gegebenen Stelle nachzuweisen, auch den größten Wert am Rand · kennst du seit Q1 · Abitur GK · baut auf: Monotonie (Einheit 1), Gleichungen lösen, Ableitungen bilden, Funktionswerte (Nr. 2–6)}

\begin{aufgabe}{Ich erkenne, ob eine Stelle gegeben ist oder gesucht wird.}
Kreuze an, ob die Stelle im Aufgabentext schon gegeben ist (nachweisen) oder erst gesucht wird (berechnen). Rechne nichts.
\begin{teile}
\teil „Bestimme die Extrempunkte des Graphen von $f$.“ \hfill \kreuz{gegeben}\kreuz{gesucht}
\teil „Zeige, dass der Graph von $f$ bei $x = 2$ einen Tiefpunkt hat.“ \hfill \kreuz{gegeben}\kreuz{gesucht}
\teil „Weise nach, dass $W(1\,|\,3)$ ein Wendepunkt des Graphen ist.“ \hfill \kreuz{gegeben}\kreuz{gesucht}
\teil „Berechne die Koordinaten des Hochpunkts.“ \hfill \kreuz{gegeben}\kreuz{gesucht}
\teil „Untersuche, ob der Graph bei $x = -1$ eine waagerechte Tangente hat.“ \hfill \kreuz{gegeben}\kreuz{gesucht}
\end{teile}
\end{aufgabe}

\begin{aufgabe}{Ich erkenne, welcher Nachweis für die Art des Punktes passt.}
Trage die passende Nummer ein. Rechne nichts.\\
(1) Vorzeichen von $f''$ \quad (2) Vorzeichenwechsel von $f'$ \quad (3) Begründung aus Abbildung oder Sachzusammenhang \quad (4) kein Nachweis verlangt
\begin{teile}
\teil Es gilt $f'(2) = 0$ und $f''(2) = -4$. \hfill Nummer: \leerfeld
\teil Es gilt $f'(0) = 0$ und $f''(0) = 0$. \hfill Nummer: \leerfeld
\teil „Bestimme die Stelle, an der $f'$ null ist. Die hinreichende Bedingung muss nicht untersucht werden.“ \hfill Nummer: \leerfeld
\teil „Bestimme den Zeitpunkt, zu dem das Wasser am höchsten steht. Begründe die Art mit dem Sachzusammenhang.“ \hfill Nummer: \leerfeld
\end{teile}
\end{aufgabe}

\verfahren{Extrempunkte berechnen}
\begin{aufgabe}{Ich kann Extrempunkte berechnen.}
Berechne die Stellen, an denen $f'(x) = 0$ ist.
\begin{teilezwei}
\tz $f(x) = x^3 - 27x$ & \tz $f(x) = x^3 + 3x^2 - 9x$ \\
\tz $f(x) = 2x^3 - 9x^2 + 12x$ & \tz $f(x) = -x^3 + 6x^2 - 9x$ \\
\anweisung{Bestimme Art und Koordinaten aller Extrempunkte.}
\tz $f(x) = x^3 - 3x^2 - 9x + 2$ & \tz $f(x) = 3x^4 - 4x^3 - 12x^2$ \\
\tz $f(x) = 3x^4 + 4x^3$ & \\
\end{teilezwei}
\begin{teile}
\teil $f(x) = x^4 - 4x^3 - 2x^2 + 12x$; bei $x = 3$ liegt eine Extremstelle.
\end{teile}
\end{aufgabe}

\begin{aufgabe}{Ich kann Extrempunkte berechnen – weiter.}
Bestimme Art und Koordinaten der Extrempunkte.
\begin{teile}
\teil $f(x) = x\cdot e^{-x}$
\anweisung{Die hinreichende Bedingung musst du nicht prüfen; begründe die Art mit dem Sachzusammenhang.}
\teil Die Konzentration eines Wirkstoffs im Blut ist $k(t) = 8t\cdot e^{-0{,}5t}$ ($t$ in Stunden, $k$ in Milligramm pro Liter). Wann ist die Konzentration am höchsten, und wie hoch ist sie?
\anweisung{Bestimme den größten Funktionswert auf dem Intervall.}
\teil $f(x) = 2x^3 - 3x^2 - 12x$ auf $[-2;\,4]$
\anweisung{Gib den Wertebereich von $f$ an.}
\teil $f(x) = x^4 - 32x$
\anweisung{Bestimme Lage und Art aller lokalen Extrempunkte.}
\teil $f(x) = (x^2 - 3)\cdot e^{-x}$ \quad (Abitur 2021 GK)
\end{teile}
\end{aufgabe}

\verfahren{Extrempunkte nachweisen}
\begin{aufgabe}{Ich kann nachweisen, dass an einer gegebenen Stelle ein Extrempunkt liegt.}
Zeige, dass $f'(x_0) = 0$ ist.
\begin{teilezwei}
\tz $f(x) = x^3 - 48x$, \ $x_0 = 4$ & \tz $f(x) = x^2 - 10x + 3$, \ $x_0 = 5$ \\
\tz $f(x) = x^4 - 4x$, \ $x_0 = 1$ & \tz $f(x) = -2x^3 + 6x$, \ $x_0 = -1$ \\
\anweisung{Weise nach, dass bei $x_0$ ein Extrempunkt liegt, und gib seine Art an.}
\tz $f(x) = x^3 + 3x^2 - 2$, \ $x_0 = -2$ & \tz $f(x) = x^4 - 5$, \ $x_0 = 0$ \\
\end{teilezwei}
\begin{teile}
\teil $f(x) = -x^3 + 12x + 3$; zeige, dass $P(2\,|\,19)$ ein Hochpunkt ist.
\anweisung{Begründe mit den gegebenen Angaben.}
\teil Es gilt $f'(1) = 0$, und der Graph von $f''$ liegt ganz unterhalb der $x$-Achse.
\anweisung{Untersuche, ob bei $x_0$ ein Extrempunkt liegt.}
\teil $f(x) = x^3 - 3x^2 + 3x$, \ $x_0 = 1$
\anweisung{Begründe ohne Rechnung.}
\teil Für $f$ gilt $f'(1) = -2$ und $f'(2) = 3$; $f'$ ist stetig und hat zwischen 1 und 2 genau eine Nullstelle. Begründe, dass $f$ zwischen 1 und 2 eine Tiefstelle hat.
\anweisung{Weise nach.}
\teil Gegeben ist $f(x) = -x^3 + 3x^2 + 1$. Weise nach, dass der Graph von $f$ den Hochpunkt $H(2\,|\,5)$ hat. (IQB 2026 grundlegend)
\end{teile}
\end{aufgabe}

\begin{aufgabe}{Ich finde den Fehler bei der Bestimmung von Extrempunkten.}
Finde den Fehler, benenne ihn und korrigiere die Rechnung.
\begin{teile}
\teil Sara untersucht $f(x) = 2x^3 - 6x + 1$:
\rechnung{f'(x) &= 6x^2 - 6 = 0 \\ x &= -1 \quad\text{oder}\quad x = 1 \\ f''(x) &= 12x, \quad f''(1) = 12 > 0}
Sara schreibt: „Bei $x = 1$ liegt ein Hochpunkt.“
\teil Ben untersucht $f(x) = (x - 2)\cdot e^{x}$:
\rechnung{f'(x) &= (x - 1)\cdot e^{x} = 0 \\ e^{x} &= 0 \quad\text{oder}\quad x - 1 = 0 \\ x &= 0 \quad\text{oder}\quad x = 1}
Ben schreibt: „Die Extremstellen liegen bei $x = 0$ und $x = 1$.“
\end{teile}
\end{aufgabe}

\begin{aufgabe}{Ich kann Aussagen über Extrempunkte begründen.}
Entscheide ohne Rechnung und begründe.
\begin{teile}
\teil Lukas sagt: „Wenn $f'(x_0) = 0$ ist, hat $f$ bei $x_0$ einen Hoch- oder Tiefpunkt.“ Stimmt das?
\teil Der Graph von $f$ hat auf $[0;\,4]$ keinen Punkt mit waagerechter Tangente. Kann $f$ auf $[0;\,4]$ trotzdem einen größten Wert haben?
\end{teile}
\end{aufgabe}

\begin{aufgabe}{Ich kann mit Extrempunkten eine Bedingung im Sachzusammenhang prüfen.}
Das Höhenprofil einer Mountainbike-Strecke mit zwei Kuppen wird für $0 \le x \le 8$ durch
\[ h(x) = -0{,}05x^4 + 0{,}8x^3 - 3{,}9x^2 + 5{,}6x + 2 \]
beschrieben ($x$ und $h$ in Längeneinheiten, 1 Längeneinheit entspricht 10 m).
\begin{teile}
\teil Bestimme die Koordinaten der beiden Kuppen.
\teil Die beiden Kuppen sollen mindestens 50 m voneinander entfernt sein. Prüfe, ob das Profil diese Bedingung erfüllt.
\end{teile}
\end{aufgabe}
